% !PRINTED 58
\weberpage{58}
%<<CONT

whence it follows that in (4) any product $x_1^{\alpha_1}x_2^{\alpha_2}\cdots$ occurs exactly
\[
 {\Pi(n)\over \Pi(\alpha_1)\Pi(\alpha_2)\cdots}
\]
times. If one further fixes that $A_{\nu_1,\nu_2,\ldots,\nu_n}$ shall not change when the indices are permuted in any way, then the notations (2) and (4) prove to be identical if, by collecting equal factors,
\[
 x_{\nu_1}x_{\nu_2}\cdots x_{\nu_n}=x_1^{\alpha_1}x_2^{\alpha_2}\cdots x_m^{\alpha_m}
\]
and
\[
 A_{\nu_1,\nu_2,\ldots,\nu_n}=A_{\alpha_1,\alpha_2,\ldots,\alpha_m}
\]
are put.

If the number of terms that occur in the function $\Phi$ [according to (2)] is denoted by $(m,n)$, then, by first counting the terms which have the factor $x_1$ and then the remaining ones, which form a homogeneous function of order $n$ of the remaining $m-1$ variables, one finds the recursion formula
\begin{equation}
 (m,n)=(m,n-1)+(m-1,n),
\tag{5}
\end{equation}
with whose aid one proves by complete induction the expression
\begin{equation}
 (m,n)={m(m+1)\cdots(m+n-1)\over 1\cdot2\cdots n}
 ={ \Pi(m+n-1)\over \Pi(n)\Pi(m-1)}
\tag{6}
\end{equation}
to be correct.

Whole homogeneous functions are also called \gesperrt{forms}. According to the number of variables one distinguishes \gesperrt{unary forms} (simple powers), \gesperrt{binary, ternary, quaternary forms}. It is the binary forms that especially interest us here; their theory is essentially identical with the theory of whole rational functions of one variable. One passes from the binary forms back to these functions by regarding one of the homogeneous variables as constant, for example by giving it the value $1$.
